\section{Rangkuman}
\label{sec:rangkuman-bab17}

Pada bab ini telah dibahas:

\begin{itemize}
    \item Perjalanan kompetensi mata kuliah ini bergerak dari
          pemahaman mesin (\CPMK{01}), ke hubungan instruksi dengan
          perangkat keras (\CPMK{02}), ke penulisan program
          (\CPMK{03}), dan berakhir pada kemampuan memperbaiki
          program (\CPMK{04}).

    \item Setiap Sub-CPMK memiliki bukti belajar yang harus
          ditunjukkan. Bukti tersebut bersifat kumulatif dan
          dipelihara sepanjang semester.

    \item Ujian Tengah Semester mengukur Sub-CPMK 1.1 sampai 2.3,
          sedangkan Ujian Akhir Semester mengukur Sub-CPMK 3.1
          sampai 4.3 secara komprehensif dengan penekanan pada
          kemampuan menulis dan memperbaiki program.

    \item Proyek akhir menawarkan empat pilihan dan dinilai dengan
          rubrik lima kriteria yang seluruhnya dapat ditelusuri ke
          CPMK.

    \item Penilaian diri dan penilaian antarteman adalah bagian
          resmi penilaian dan harus disertai bukti.

    \item Nilai rubrik tanpa alasan hanya menghasilkan angka; rubrik
          dengan alasan menghasilkan umpan balik yang dapat
          ditindaklanjuti.

    \item Tindak lanjut dapat diarahkan ke arsitektur komputer,
          pemrograman sistem, analisis keamanan, atau sistem
          tertanam, sesuai minat mahasiswa.
\end{itemize}

\begin{checklistbox}
\textbf{Checklist Kompetensi Akhir Mata Kuliah}

Tandai butir yang benar-benar dapat Anda buktikan dengan berkas
atau catatan.

\begin{itemize}
    \item[\checkbox] Saya dapat menjelaskan arsitektur 8086, register,
          dan segmentasi memori, serta menunjukkan diagram dan
          perhitungan yang mendukungnya.
    \item[\checkbox] Saya dapat membaca format instruksi mesin,
          menentukan mode pengalamatan, dan menjelaskan alur translasi
          program.
    \item[\checkbox] Saya dapat menulis program assembly lengkap yang
          memuat aritmatika, logika, struktur kontrol, dan prosedur
          berparameter.
    \item[\checkbox] Saya dapat menggunakan interupsi BIOS/DOS untuk
          operasi masukan dan keluaran.
    \item[\checkbox] Saya dapat menelusuri program langkah demi
          langkah, membaca isi memori dan stack, serta menemukan
          penyebab kesalahan.
    \item[\checkbox] Saya dapat membaca hasil \eng{disassembly} dan
          mengusulkan perbaikan yang terukur.
    \item[\checkbox] Portofolio, jurnal praktikum, dan laporan proyek
          akhir saya lengkap.
\end{itemize}

\vspace{0.4em}
Seluruh butir bertanda berarti seluruh CPMK mata kuliah ini telah
tercapai.
\end{checklistbox}

\begin{infobox}
\textbf{Rujukan bab.} Untuk memperluas pemahaman mengenai organisasi
komputer yang menjadi dasar mata kuliah ini, lihat
\cite{stallings2016} dan \cite{tanenbaum2013}. Untuk jalur
pemrograman sistem dan analisis level rendah, lihat
\cite{bartlett2004}, \cite{bartlett2023}, dan
\cite{wikibooksdisasm}. Untuk rujukan teknis resmi keluarga x86,
lihat \cite{intel2024}.
\end{infobox}
